% 数学ノート集・共通テンプレート v2（積分論2026年度第1部の版面を参考）
% A5判。章内で定義・定理・例・注意等の番号を共通化。
% コンパイル： platex -interaction=nonstopmode math_note_template_v2.tex (2回)
%              dvipdfmx math_note_template_v2.dvi
\documentclass[a5paper,10pt,onecolumn]{jsarticle}
\usepackage[dvipdfmx,top=22mm,bottom=20mm,inner=19mm,outer=18mm,headheight=13pt,headsep=6mm]{geometry}
\usepackage{amsmath,amssymb,amsthm,mathtools}
\usepackage{enumitem,fancyhdr,titlesec}
\usepackage[dvipdfmx,bookmarks=true,colorlinks=false]{hyperref}
\usepackage{pxjahyper}

% ----- 差し替える文書情報 -----
\newcommand{\NoteTitle}{Frobeniusの定理}
\newcommand{\NoteSeries}{数学ノート集}

% ----- 文字・組版 -----
% 参考PDFの質素な明朝＋Computer Modernの印象を引き継ぎながら，
% A5判で窮屈にならないよう本文の行間を少し確保する．
\linespread{1.13}
\setlength{\parindent}{1zw}
\setlength{\parskip}{0pt}
\setlength{\abovedisplayskip}{6pt plus 2pt minus 2pt}
\setlength{\belowdisplayskip}{6pt plus 2pt minus 2pt}
\setlength{\abovedisplayshortskip}{3pt plus 2pt minus 1pt}
\setlength{\belowdisplayshortskip}{4pt plus 2pt minus 1pt}
\setlist[enumerate]{topsep=3pt,itemsep=1pt,parsep=0pt,leftmargin=2.4zw}
\titleformat{\section}{\normalfont\normalsize\bfseries}{\thesection}{0.8zw}{}
\titlespacing*{\section}{0pt}{2.0ex plus .4ex}{.9ex}
\titleformat{\subsection}{\normalfont\normalsize\bfseries}{\thesubsection}{0.7zw}{}
\titlespacing*{\subsection}{0pt}{1.5ex plus .3ex}{.6ex}
\raggedbottom

% ----- 柱とページ番号：ページ上部、罫線なし -----
\fancypagestyle{mathnote}{%
  \fancyhf{}%
  \fancyhead[L]{\footnotesize \NoteSeries\ （\NoteTitle）}%
  \fancyhead[R]{\footnotesize \thepage}%
  \renewcommand{\headrulewidth}{0pt}%
  \renewcommand{\footrulewidth}{0pt}%
}
\pagestyle{mathnote}

% ----- 節内通し番号の定義・定理群 -----
\newtheoremstyle{mathnotetheorem}
 {4.5pt}{4.5pt}{\normalfont}{}%
 {\bfseries}{}{0.65zw}%
 {\thmname{#1}\thmnumber{ #2}\thmnote{（#3）}}
\theoremstyle{mathnotetheorem}
\newtheorem{theorem}{定理}[section]
\newtheorem{definition}[theorem]{定義}
\newtheorem{lemma}[theorem]{補題}
\newtheorem{proposition}[theorem]{命題}
\newtheorem{corollary}[theorem]{系}
\newtheorem{example}[theorem]{例}
\newtheorem{remark}[theorem]{注意}
\renewcommand{\proofname}{証明}
\renewcommand{\qedsymbol}{$\square$}
\makeatletter
\renewenvironment{proof}[1][\proofname]{%
  \par\pushQED{\qed}\normalfont\topsep4\p@\@plus2\p@\relax
  \trivlist\item[\hskip\labelsep\bfseries #1]\ignorespaces
}{\popQED\endtrivlist\@endpefalse}
\makeatother

% ----- 使うノートで適宜拡張する記号 -----
\newcommand{\R}{\mathbb{R}}
\newcommand{\norm}[1]{\lVert #1\rVert}
\newcommand{\abs}[1]{\lvert #1\rvert}

\begin{document}
\section*{読了後にできること}
本稿を読むための前提は，分布の定義，Lie括弧，Frobeniusの定理．
本稿の具体的な到達課題は，与えられたベクトル場族の可積分性をLie括弧で判定する．


\section{分布と積分多様体}
\begin{definition}
$M$ 上の階数 $k$ の滑らかな分布 $\mathcal D$ とは，
各点の接空間の $k$ 次元部分空間 $\mathcal D_p\subset T_pM$
の滑らかな局所フレームによって定まる部分束をいう．
滑らかなベクトル場 $X,Y$ が
$\mathcal D$ に値を取るとき常に
$[X,Y]$ も $\mathcal D$ に値を取るなら
$\mathcal D$ は対合的という．
\end{definition}
\begin{definition}
埋め込み部分多様体 $N\subset M$ が分布 $\mathcal D$ の
積分多様体とは，任意の $p\in N$ で
$T_pN=\mathcal D_p$ となることをいう．
\end{definition}
\begin{proposition}[積分多様体から対合性]
$\mathcal D$ が各点を通る局所積分多様体を持つなら
$\mathcal D$ は対合的である．
\end{proposition}
\begin{proof}
局所積分多様体 $N$ を固定すると，
$\mathcal D$ に値を取るベクトル場 $X,Y$ の制限は
$N$ に接する．
$N$ を定義する局所方程式 $f^\alpha=0$ に対し
$Xf^\alpha=Yf^\alpha=0$ が $N$ 上で成立する．
従って $[X,Y]f^\alpha=X(Yf^\alpha)-Y(Xf^\alpha)=0$
が $N$ 上で成立し，Lie括弧も $N$ に接する．
各点で $\mathcal D_p=T_pN$ なので対合性を得る．
\end{proof}
\section{正規化した局所フレーム}
\begin{lemma}[正規化フレーム]\label{lem:normalframe}
階数 $k$ の分布 $\mathcal D$ に対し，
各点の近傍で局所座標 $x^1,\dots,x^n$ と
$\mathcal D$ のフレーム
\[
 X_i=\frac{\partial}{\partial x^i}
 +\sum_{\alpha=k+1}^n a_i^\alpha(x)
 \frac{\partial}{\partial x^\alpha},
 \qquad 1\le i\le k
\]
を取れる．
\end{lemma}
\begin{proof}
点 $p$ における $\mathcal D_p$ の基底を
$\mathbb R^n$ の最初の $k$ 個の座標方向となるように
線形座標変換する．
局所フレーム $Y_i$ の第1から第$k$成分で作る
$k\times k$ 行列は $p$ で可逆なので近傍でも可逆である．
その逆行列を係数として $Y_i$ の線形結合を取ると，
最初の $k$ 成分が単位行列となり表示式を得る．
\end{proof}
\begin{lemma}[対合性と可換性]\label{lem:commuting}
補題\ref{lem:normalframe}のフレームについて，
$\mathcal D$ が対合的なら
$[X_i,X_j]=0$（全ての $i,j$）である．
\end{lemma}
\begin{proof}
$[X_i,X_j]$ の初めの $k$ 座標成分は
$X_i(\delta_j^\ell)-X_j(\delta_i^\ell)=0$ である．
対合性から $[X_i,X_j]=\sum_\ell b^\ell X_\ell$ と書ける．
初めの $k$ 成分を比較すると $b^\ell=0$ であり，
Lie括弧は零となる．
\end{proof}
\section{Frobeniusの定理}
\begin{theorem}[Frobeniusの可積分性定理]\label{thm:frobenius}
滑らかな階数一定の分布 $\mathcal D$ に対し，
次は同値である．
\begin{enumerate}
\item $\mathcal D$ は対合的である．
\item 各点の近傍で
 $\mathcal D=\operatorname{span}
 \{\partial/\partial y^1,\ldots,\partial/\partial y^k\}$
 となる局所座標が存在する．
\item 各点を通る局所積分多様体が存在する．
\end{enumerate}
\end{theorem}
\begin{proof}
(2) なら $y^{k+1},\dots,y^n$ を定数とする
座標切片が積分多様体となるので (3) が従う．
(3) から (1) は前命題で示した．
(1) から (2) を示す．
補題\ref{lem:normalframe}の正規化フレーム $X_i$ を取り，
補題\ref{lem:commuting}によって互いに可換である．
各 $X_i$ の局所フローを $\Phi_i^t$ とする．
Lie括弧が零なのでフローは定義される範囲で可換である．
これは
\[
 \frac d{dt}(\Phi_i^{-t})_*X_j
 =(\Phi_i^{-t})_*[X_i,X_j]=0
\]
を用いて $(\Phi_i^t)_*X_j=X_j$ を得ること，
および同じ初期値からの積分曲線の一意性によって従う．
$p$ の座標で，横断切片
$\sigma(z)=(0,\dots,0,z^{k+1},\dots,z^n)$ を取り
\[
 F(t^1,\dots,t^k,z)=
 \Phi_1^{t^1}\circ\cdots\circ\Phi_k^{t^k}(\sigma(z))
\]
を定義する．
$DF_{(0,z_0)}$ の初めの $k$ 列は $X_i(p)$，
残りの列は $\partial_\alpha(p)$（$\alpha>k$）である．
正規化によりこれらは接空間の基底なので
逆写像定理から $F$ は局所微分同相である．
フローの可換性と生成性によって
$F_*(\partial/\partial t^i)=X_i$ であり，
$F$ の逆座標 $(t,z)$ では
$\mathcal D=\operatorname{span}\{\partial/\partial t^i\}$
となる．よって (2) が成立する．
\end{proof}
\section{1形式による記述}
\begin{proposition}[余次元1のFrobenius条件]
$\alpha$ を零にならない滑らかな1形式とし，
$\mathcal D=\ker\alpha$ とする．
このとき $\mathcal D$ が対合的であることと
\[
 \alpha\wedge d\alpha=0
\]
は同値である．
\end{proposition}
\begin{proof}
$X,Y\in\Gamma(\mathcal D)$ について
外微分の定義から
$d\alpha(X,Y)=X(\alpha(Y))-Y(\alpha(X))-\alpha([X,Y])
=-\alpha([X,Y])$．
従って対合性は $d\alpha$ が
$\mathcal D\times\mathcal D$ 上零であることと同値である．
局所的に $\alpha=e^1$ となる余接フレームを取り
$d\alpha=\sum_{i<j}b_{ij}e^i\wedge e^j$ と書く．
$\alpha\wedge d\alpha=0$ は $b_{ij}=0$（$i,j>1$）
と同値であり，これは
$d\alpha|_{\mathcal D\times\mathcal D}=0$ と同値である．
\end{proof}
\begin{example}[非可積分な接平面場]
$\mathbb R^3$ 上の $\alpha=dz-y\,dx$ と
$\mathcal D=\ker\alpha$ を考える．これは対合的でない．
\end{example}
\begin{proof}
$d\alpha=-dy\wedge dx=dx\wedge dy$ なので
$\alpha\wedge d\alpha=dz\wedge dx\wedge dy\ne0$．
前命題から対合的でなく，積分曲面は存在しない．
\end{proof}

\section{到達課題の完全解答}
\noindent\textbf{接触分布の非可積分性を確かめる．}
\begin{proposition}[非可積分な二次元分布]
$\mathbb R^3$ のベクトル場
\[
X=\partial_x+y\partial_z,\qquad
Y=\partial_y
\]
の張る階数2の分布 $D$ は可積分でない．
\end{proposition}
\begin{proof}
Lie括弧は
\[
[X,Y]=-\partial_z
\]
である．
もし $-\partial_z=aX+bY$ と書ければ
$\partial_x$ 成分から $a=0$，
$\partial_y$ 成分から $b=0$ となり矛盾する．
よって $[X,Y]\notin D$．
Frobeniusの定理の必要条件を満たさないため
$D$ は可積分でない．
\end{proof}
\begin{proposition}[可積分な比較例]
$X=\partial_x$，$Y=\partial_y$ の張る分布は可積分であり，
積分多様体は $z=c$ の平面である．
\end{proposition}
\begin{proof}
$[X,Y]=0$ だから括弧で閉じている．
また $z=c$ 上では接空間が
$\operatorname{span}\{\partial_x,\partial_y\}$
と一致するので，実際に積分多様体を構成できる．
\end{proof}
\noindent\textbf{到達確認．}
$X=\partial_x+x\partial_z$，
$Y=\partial_y$ の分布の可積分性を判定せよ．
\par\noindent\textbf{解答．}
$[X,Y]=0$ だからFrobeniusの定理より可積分である．
実際 $z-x^2/2=c$ の曲面が積分多様体となる．
\end{document}
